% Post-CP17 refinement; full root and obstruction/crystal audits passed.
\subsection{The logarithm equality boundary and square finite determinants}

\begin{lemma}[The closed logarithm boundary]
\label{lem:B-RD50-boundary}
\status{Proved; complete root and obstruction/crystal proof, original-source, and exact-code audits passed.}
Let $R$ be a complete discrete valuation ring of mixed characteristic
$(0,p)$ with perfect residue field, and put $a=v(p)$ for its normalized
valuation. Let $f:F\to G$ be a homomorphism of finite-height formal
groups with no nonzero finite flat subgroup in its kernel and with
finite special-fibre kernel. Suppose $G_k$ is isomorphic to a product
of one-dimensional height-two formal groups. If
\[
                              a\leq p^2-1,
                                                        \tag{B-RD50.1}
\]
then relative Frobenius annihilates $\ker f_k$.
\end{lemma}
\begin{proof}
We keep the actual Cartier-module argument of
Lemma~\ref{lem:B-RD34-kernel}, and verify its coefficient bound at
an equality previously left unused. The normalized invariant
differential matrix gives
\[
                  U([p]X)J_{[p]}(X)=pU(X).
\]
Every monomial coefficient with an exponent not divisible by $p$
is consequently divisible by $p$. In particular all nonlinear
terms of total degree less than $p$ are divisible by $p$.
The product-height-two special-fibre condition says that every
coefficient of total degree $p$ through $p^2-1$ lies in the
maximal ideal. It is invariant under integral changes of formal
coordinates, exactly as proved in Lemma~\ref{lem:B-RD34-log}.
Thus for any augmentation vector $h\in tR[[t]]^{\dim G}$,
\[
        [p]_G(\pi h)\in\pi^qR[[t]]^{\dim G},\qquad
        q=\min\{a+1,p+1,p^2\}\geq a/(p-1).
                                                        \tag{B-RD50.2}
\]
The last inequality follows from (B-RD50.1).

For a coefficient $c_\alpha$ of the normalized logarithm, of
positive total degree $n=|\alpha|$, integration of the invariant
differentials gives
\[
 v(c_\alpha)\geq-a\lfloor\log_p n\rfloor,
 \qquad \lfloor\log_p n\rfloor\leq\frac{n-1}{p-1}.
\]
It follows that if all coefficients of an augmentation vector
$z$ have valuation at least $q\geq a/(p-1)$, then every
coefficient of $\log_G(z)$ has valuation at least $q$.
Equality in the bound causes no problem: formal substitution
in the augmentation variable $t$ involves finitely many terms
in each coefficient. No analytic-convergence assertion at the
boundary is used. Apply this to (B-RD50.2) to obtain
\[
                      p\log_G(\pi h)\in\pi^qR[[t]]^{\dim G}.
                                                        \tag{B-RD50.3}
\]

For precision, put $J=\pi^{\lceil a/(p-1)\rceil}R$ and use
Bondarko's modules $D_G,\mathcal C_G$ and his bounded-series ring
$\mathcal R=R[[\Delta]][1/p]$. His exact identities are
\[
 \mathcal C_G=\Omega\pi D_{G_\pi},\quad
 J[[\Delta]]^{\dim G}\subset D_G,\quad
 E(G_k)=D_G/(D_G\cap\mathcal R^{\dim G}).
\]
The ideal in his Proposition~9 is
$\mathfrak m^{-[a/(1-p)]}=J$, including the integer equality
case. The precise source is
\cite[Propositions~9--12 and Lemma~3, pp.~674--676;
Proposition~16 and its full proof, p.~681]{Bondarko2006Generic}.
The integral Witt-coefficient action of $\Omega$ preserves each
coefficient-valuation ideal. Equation (B-RD50.3) therefore gives
$p\mathcal C_G\subset J\{\{\Delta\}\}^{\dim G}$.
The bounded-projection identity in Proposition~11(2), together
with $\mathbf fz=p(z-z(0))\Delta^{-1}$ and $pR\subset J$,
then gives
\[
 z\in D_G,\ Az\in\mathcal R^b
       \quad\Longrightarrow\quad
                 \mathbf fAz\in J[[\Delta]]^b
                                                        \tag{B-RD50.4}
\]
for every integral matrix $A$. This is the exact replacement
for (B-RD34.8), with one Cartier Frobenius.

The remaining saturation argument is unchanged, but its target
is explicitly $G$. Put $I'=df(0)(\operatorname{Frac}(R)^{\dim F})$,
choose an integral complement to $I'\cap R^{\dim G}$, and let
$\operatorname{pr}$ be the associated projection. Absence of a
finite flat kernel gives
\[
 df(0)D_F=D_G\cap I'[[\Delta]].
\]
For a class in the saturation of $E(F_k)$ inside $E(G_k)$,
choose $z\in D_G$. Then $\operatorname{pr}z$ is bounded.
By (B-RD50.4), $c=\mathbf f\operatorname{pr}z$ lies in
$J[[\Delta]]$, hence in $D_G$, and is bounded. Thus
$\mathbf fz-c\in df(0)D_F$, while $c$ reduces to zero.
The saturation quotient is annihilated by $\mathbf f$.
By \cite[Theorem~2(8), equation~(6), p.~668]{Bondarko2006Generic},
this is the Dieudonne module of $\ker f_k$, and $\mathbf f$
corresponds to relative Frobenius. This proves the assertion.
\end{proof}

\begin{proposition}[The square-determinant tame comparison]
\label{prop:B-RD50-square}
\status{Proved; complete root and obstruction/crystal proof, original-source, and exact-code audits passed.}
Let $p\geq5$ and let $K$ be complete discretely valued of mixed
characteristic $(0,p)$, with perfect residue field,
$e_K=p-1$, and $\mu_p\subset K$.
Suppose two height-two groups satisfying the manuscript monodromy
condition have isomorphic $p^2$-torsion and different requested
characters. After completed maximal unramified base change,
Theorem~\ref{thm:B-RD49} makes both groups potentially
supersingular, and their determinants have the common form
$\chi_p\psi$, with $\psi:I_K\to\mu_{p-1}$ finite.
Then $\psi$ has no square root in
$\operatorname{Hom}_{\mathrm{cont}}(I_K,\mu_{p-1})$.
In particular the joint faithful tame degree cannot equal $p+1$.
\end{proposition}
\begin{proof}
Assume instead that $\psi=\epsilon^2$ for a finite
$\mu_{p-1}$-valued character. Twist both groups by
$\epsilon^{-1}$. This preserves the torsion isomorphism and
multiplies both special characters by the same nonzero scalar
character, so it preserves inequality. Both new determinants
are cyclotomic. Their special finite actions have reduced norm
one, by the determinant comparison already used in
Proposition~\ref{prop:B-RD32}. The height-two division algebra
calculation in Lemma~\ref{lem:B-RD49-inertia} shows that their
finite images are cyclic of orders dividing $p+1$.
Indeed they are prime-to-$p$ roots of unity in an unramified
quadratic field, and the reduced norm is the $(p+1)$st power.
Their common faithful good field $L$ is therefore tame cyclic
of degree $m\mid p+1$, and $a=e_L=(p-1)m\leq p^2-1$.

Let $\mathcal M,\mathcal N$ be their good formal models and
let $\tau$ generate $\operatorname{Gal}(L/K)$.
Different norm-one characters mean that the two unordered
special eigenvalue sets $\{\zeta,\zeta^{-1}\}$ are disjoint.
Their minimal polynomials over $\mathbf Z_p$ have relatively
prime reductions: the roots are distinct Teichmuller roots
of unity, with their $p$th-power conjugate equal to their
inverse. The Chinese remainder theorem consequently gives
an integral polynomial in $\tau$ which is the projection
onto one factor of $M_0\times N_0$.

Let $H$ be the finite flat closure of the graph of the given
isomorphism in $\mathcal M[p^2]\times\mathcal N[p^2]$.
The perfect alternating pairings needed here are actual
integral BT pairings: cyclotomic determinants give them on
the full generic Tate lattices, and Tate's theorem extends
both the duality map and its inverse over $\mathcal O_L$.
This is exactly the argument in Proposition~\ref{prop:B-RD33},
using \cite[Theorem~4 and Corollaries~1--2, p.~181]
{Tate1967PDivisible}. Scale one pairing by the unit multiplier
of the torsion isomorphism. The graph closure is then its own
orthogonal, including after reduction, by flatness and equal
rank as in that proposition.

The integral tame projector preserves $H_k$, so
$H_k=J_M\times J_N$. A one-dimensional connected formal group
over an algebraically closed field has a unique subgroup of
rank $p^r$, namely its relative $F^r$-kernel. Orthogonality in
the perfectly paired $p^2$-torsion makes both exponents equal
to two. Thus
\[
                    H_k=M_0[p]\times N_0[p].            \tag{B-RD50.5}
\]
Set $Q=(\mathcal M\times\mathcal N)/H$. Its special fibre is
$M_0^{(p^2)}\times N_0^{(p^2)}$, a product of height-two groups.
The homomorphism
\[
 \mathcal M\times\mathcal N\longrightarrow Q\times\mathcal M,
              \qquad z\longmapsto(qz,\operatorname{pr}_{\mathcal M}z)
\]
has no nonzero finite flat subgroup in its kernel, because its
generic graph meets $0\times N$ trivially. Its special kernel
is $0\times N_0[p]=0\times\ker F_{N_0}^2$.
Lemma~\ref{lem:B-RD50-boundary} says that relative Frobenius
annihilates that kernel. Its Frobenius image instead has
rank $p$, a contradiction.

Finally, if the original joint degree were $p+1$, both original
special norm characters would already be trivial, so the same
argument would apply without twisting.
\end{proof}

\begin{corollary}[An explicit normalized remaining character class]
\label{cor:B-RD50-remaining}
\status{Proved; complete root and obstruction/crystal proof, original-source, and exact-code audits passed.}
In the remaining $p^2$ test of Theorem~\ref{thm:B-RD49}, the
common finite determinant must have nonsquare value on a tame
generator. After a simultaneous finite height-one twist one
may require that this value be primitive of order $p-1$.
In that normalization the joint faithful tame degree has the form
\[
          m=\frac{p^2-1}{d},\qquad
          d\mid\frac{p+1}{2^{v_2(p+1)}}.
                                                        \tag{B-RD50.6}
\]
For $p=5$ this leaves $m=8,24$; for $p=7$ it leaves only
$m=48$. These are necessary conditions, with no assertion that
any remaining pair exists. The least-index interval remains
$\{1,2\}$.
\end{corollary}
\begin{proof}
The tame inertia group is procyclic. A character into the cyclic
group $\mu_{p-1}$ has a square root precisely when its value on
a generator is a square. Choose a character $\eta$ with primitive
value. For nonsquare $\psi$, the character $\psi\eta^{-1}$ is
a square, so a simultaneous twist makes the determinant character
$\eta$.

A special tame eigenvalue is a Teichmuller root in
$\mu_{p^2-1}$, say $\zeta^c$ for a fixed primitive root.
Its reduced norm is $(\zeta^{p+1})^c$. Primitivity of the finite
determinant implies $\gcd(c,p-1)=1$ (the inverse norm convention,
if used, has the same consequence). Therefore $c$ is odd and
$\gcd(c,p^2-1)$ is an odd divisor of $p+1$, since
$\gcd(p-1,p+1)=2$. Each individual faithful order is consequently
$(p^2-1)/d_i$ with $d_i$ an odd divisor of $p+1$. Their joint
order is the least common multiple, obtained with
$d=\gcd(d_1,d_2)$, proving (B-RD50.6).
\end{proof}

\begin{remark}[Boundary scope and exact certificate]
\label{rem:B-RD50-scope}
\status{Proved under stated hypotheses.}
The new assertion uses the closed coefficientwise logarithm
bound. It neither assumes analytic convergence at equality nor
transfers Hasse-coefficient thickness from $\mathcal M,\mathcal N$
to the quotient $Q$. Only the product shape of $Q_k$ is used.
The earlier strict logarithm lemmas remain correct within their
stated ranges; this argument supplies an additional boundary case.
The retained exact program
\path{code/reduction_descent/verify_tame_norm_square_boundary.py}
checks 8,243,000 coefficientwise logarithm inequalities, 4,128
coprime tame-spectrum pairs, and 2,908 primitive-norm exponents.
The author run passed in 1.6194844739948167 seconds; the
independent root and obstruction/crystal runs passed in
1.491409575995931 and 1.5930737649978255 seconds, respectively.
Both independent runs preserved the author output. The actual
Cartier-module proof and original published Bondarko passages
were independently reread in full. These finite checks do not
assert the existence of any remaining torsion-congruent pair.
\end{remark}
